% Options for packages loaded elsewhere
\PassOptionsToPackage{unicode}{hyperref}
\PassOptionsToPackage{hyphens}{url}
\PassOptionsToPackage{dvipsnames,svgnames,x11names}{xcolor}
%
\documentclass[
]{scrartcl}

\usepackage{amsmath,amssymb}
\usepackage{iftex}
\ifPDFTeX
  \usepackage[T1]{fontenc}
  \usepackage[utf8]{inputenc}
  \usepackage{textcomp} % provide euro and other symbols
\else % if luatex or xetex
  \usepackage{unicode-math}
  \defaultfontfeatures{Scale=MatchLowercase}
  \defaultfontfeatures[\rmfamily]{Ligatures=TeX,Scale=1}
\fi
\usepackage{lmodern}
\ifPDFTeX\else  
    % xetex/luatex font selection
\fi
% Use upquote if available, for straight quotes in verbatim environments
\IfFileExists{upquote.sty}{\usepackage{upquote}}{}
\IfFileExists{microtype.sty}{% use microtype if available
  \usepackage[]{microtype}
  \UseMicrotypeSet[protrusion]{basicmath} % disable protrusion for tt fonts
}{}
\makeatletter
\@ifundefined{KOMAClassName}{% if non-KOMA class
  \IfFileExists{parskip.sty}{%
    \usepackage{parskip}
  }{% else
    \setlength{\parindent}{0pt}
    \setlength{\parskip}{6pt plus 2pt minus 1pt}}
}{% if KOMA class
  \KOMAoptions{parskip=half}}
\makeatother
\usepackage{xcolor}
\setlength{\emergencystretch}{3em} % prevent overfull lines
\setcounter{secnumdepth}{5}
% Make \paragraph and \subparagraph free-standing
\ifx\paragraph\undefined\else
  \let\oldparagraph\paragraph
  \renewcommand{\paragraph}[1]{\oldparagraph{#1}\mbox{}}
\fi
\ifx\subparagraph\undefined\else
  \let\oldsubparagraph\subparagraph
  \renewcommand{\subparagraph}[1]{\oldsubparagraph{#1}\mbox{}}
\fi

% Lower the position of page numbers
\usepackage{fancyhdr}
\fancyhf{}
\fancyfoot[C]{\thepage}
\setlength{\footskip}{100pt} % Increase this value to lower the page numbers


\usepackage{color}
\usepackage{fancyvrb}
\newcommand{\VerbBar}{|}
\newcommand{\VERB}{\Verb[commandchars=\\\{\}]}
\DefineVerbatimEnvironment{Highlighting}{Verbatim}{commandchars=\\\{\}}
% Add ',fontsize=\small' for more characters per line
\usepackage{framed}
\definecolor{shadecolor}{RGB}{241,243,245}
\newenvironment{Shaded}{\begin{snugshade}}{\end{snugshade}}
\newcommand{\AlertTok}[1]{\textcolor[rgb]{0.68,0.00,0.00}{#1}}
\newcommand{\AnnotationTok}[1]{\textcolor[rgb]{0.37,0.37,0.37}{#1}}
\newcommand{\AttributeTok}[1]{\textcolor[rgb]{0.40,0.45,0.13}{#1}}
\newcommand{\BaseNTok}[1]{\textcolor[rgb]{0.68,0.00,0.00}{#1}}
\newcommand{\BuiltInTok}[1]{\textcolor[rgb]{0.00,0.23,0.31}{#1}}
\newcommand{\CharTok}[1]{\textcolor[rgb]{0.13,0.47,0.30}{#1}}
\newcommand{\CommentTok}[1]{\textcolor[rgb]{0.37,0.37,0.37}{#1}}
\newcommand{\CommentVarTok}[1]{\textcolor[rgb]{0.37,0.37,0.37}{\textit{#1}}}
\newcommand{\ConstantTok}[1]{\textcolor[rgb]{0.56,0.35,0.01}{#1}}
\newcommand{\ControlFlowTok}[1]{\textcolor[rgb]{0.00,0.23,0.31}{\textbf{#1}}}
\newcommand{\DataTypeTok}[1]{\textcolor[rgb]{0.68,0.00,0.00}{#1}}
\newcommand{\DecValTok}[1]{\textcolor[rgb]{0.68,0.00,0.00}{#1}}
\newcommand{\DocumentationTok}[1]{\textcolor[rgb]{0.37,0.37,0.37}{\textit{#1}}}
\newcommand{\ErrorTok}[1]{\textcolor[rgb]{0.68,0.00,0.00}{#1}}
\newcommand{\ExtensionTok}[1]{\textcolor[rgb]{0.00,0.23,0.31}{#1}}
\newcommand{\FloatTok}[1]{\textcolor[rgb]{0.68,0.00,0.00}{#1}}
\newcommand{\FunctionTok}[1]{\textcolor[rgb]{0.28,0.35,0.67}{#1}}
\newcommand{\ImportTok}[1]{\textcolor[rgb]{0.00,0.46,0.62}{#1}}
\newcommand{\InformationTok}[1]{\textcolor[rgb]{0.37,0.37,0.37}{#1}}
\newcommand{\KeywordTok}[1]{\textcolor[rgb]{0.00,0.23,0.31}{\textbf{#1}}}
\newcommand{\NormalTok}[1]{\textcolor[rgb]{0.00,0.23,0.31}{#1}}
\newcommand{\OperatorTok}[1]{\textcolor[rgb]{0.37,0.37,0.37}{#1}}
\newcommand{\OtherTok}[1]{\textcolor[rgb]{0.00,0.23,0.31}{#1}}
\newcommand{\PreprocessorTok}[1]{\textcolor[rgb]{0.68,0.00,0.00}{#1}}
\newcommand{\RegionMarkerTok}[1]{\textcolor[rgb]{0.00,0.23,0.31}{#1}}
\newcommand{\SpecialCharTok}[1]{\textcolor[rgb]{0.37,0.37,0.37}{#1}}
\newcommand{\SpecialStringTok}[1]{\textcolor[rgb]{0.13,0.47,0.30}{#1}}
\newcommand{\StringTok}[1]{\textcolor[rgb]{0.13,0.47,0.30}{#1}}
\newcommand{\VariableTok}[1]{\textcolor[rgb]{0.07,0.07,0.07}{#1}}
\newcommand{\VerbatimStringTok}[1]{\textcolor[rgb]{0.13,0.47,0.30}{#1}}
\newcommand{\WarningTok}[1]{\textcolor[rgb]{0.37,0.37,0.37}{\textit{#1}}}

\providecommand{\tightlist}{%
  \setlength{\itemsep}{0pt}\setlength{\parskip}{0pt}}\usepackage{longtable,booktabs,array}
\usepackage{calc} % for calculating minipage widths
% Correct order of tables after \paragraph or \subparagraph
\usepackage{etoolbox}
\makeatletter
\patchcmd\longtable{\par}{\if@noskipsec\mbox{}\fi\par}{}{}
\makeatother
% Allow footnotes in longtable head/foot
\IfFileExists{footnotehyper.sty}{\usepackage{footnotehyper}}{\usepackage{footnote}}
\makesavenoteenv{longtable}
\usepackage{graphicx}
\makeatletter
\newsavebox\pandoc@box
\newcommand*\pandocbounded[1]{% scales image to fit in text height/width
  \sbox\pandoc@box{#1}%
  \Gscale@div\@tempa{\textheight}{\dimexpr\ht\pandoc@box+\dp\pandoc@box\relax}%
  \Gscale@div\@tempb{\linewidth}{\wd\pandoc@box}%
  \ifdim\@tempb\p@<\@tempa\p@\let\@tempa\@tempb\fi% select the smaller of both
  \ifdim\@tempa\p@<\p@\scalebox{\@tempa}{\usebox\pandoc@box}%
  \else\usebox{\pandoc@box}%
  \fi%
}
% Set default figure placement to htbp
\def\fps@figure{htbp}
\makeatother

% Change font size
\usepackage[font={footnotesize, sf}, labelfont=bf, margin=0.05\linewidth]{caption}
\renewcommand{\familydefault}{\sfdefault}
% Bold the 'Figure #' in the caption and separate it from the title/caption
% with a period
% Captions will be left justified
\usepackage[
	aboveskip=1pt,
	labelfont=bf,
	labelsep=period,
	justification=raggedright,
	singlelinecheck=off
]{caption}
\renewcommand{\figurename}{Fig}

% Add numbered lines
\usepackage{lineno}
% \linenumbers

% Package to include multiple title pages
% This will allow me to add a tile to the main text and to the SI
\usepackage{titling}

% Define command to begin the supplementary section
\newcommand{\beginsupplement}{
\setcounter{page}{1} % Restart page counter
\setcounter{section}{0} % Restart section counter
% \renewcommand{\thesection}{\Alph{section}}%
\renewcommand{\thesection}{}
\renewcommand{\thesubsection}{\Alph{subsection}}
\setcounter{table}{0} % Restart table counter
\renewcommand{\thetable}{\Alph{table}}
\setcounter{figure}{0} % Restart figure counter
\renewcommand{\thefigure}{\Alph{figure}}%
\setcounter{equation}{0} % Restart equation counter
\renewcommand{\theequation}{S\arabic{equation}}%
}

% Add author affiliations
\usepackage{authblk}

% Allow to define personalized colors
\usepackage{xcolor}
\makeatletter
\@ifpackageloaded{caption}{}{\usepackage{caption}}
\AtBeginDocument{%
\ifdefined\contentsname
  \renewcommand*\contentsname{Table of contents}
\else
  \newcommand\contentsname{Table of contents}
\fi
\ifdefined\listfigurename
  \renewcommand*\listfigurename{List of Figures}
\else
  \newcommand\listfigurename{List of Figures}
\fi
\ifdefined\listtablename
  \renewcommand*\listtablename{List of Tables}
\else
  \newcommand\listtablename{List of Tables}
\fi
\ifdefined\figurename
  \renewcommand*\figurename{Figure}
\else
  \newcommand\figurename{Figure}
\fi
\ifdefined\tablename
  \renewcommand*\tablename{Table}
\else
  \newcommand\tablename{Table}
\fi
}
\@ifpackageloaded{float}{}{\usepackage{float}}
\floatstyle{ruled}
\@ifundefined{c@chapter}{\newfloat{codelisting}{h}{lop}}{\newfloat{codelisting}{h}{lop}[chapter]}
\floatname{codelisting}{Listing}
\newcommand*\listoflistings{\listof{codelisting}{List of Listings}}
\makeatother
\makeatletter
\makeatother
\makeatletter
\@ifpackageloaded{caption}{}{\usepackage{caption}}
\@ifpackageloaded{subcaption}{}{\usepackage{subcaption}}
\makeatother
\ifLuaTeX
  \usepackage{selnolig}  % disable illegal ligatures
\fi

%%%%%%%%%%%%%%%%%%%%%%%%%%%%%%%%%%%%%%%%%%%%%%%%%%%%%%%%%%%%%%%%%%%%%%%%%%%%%%%%
\usepackage[
    style=vancouver,
    sorting=none,				% Do not sort bibliography
	url=false, 					% Do not show url in reference
	doi=false, 					% Do not show doi in reference
	isbn=false, 				% Do not show isbn link in reference
	eprint=false, 			    % Do not show eprint link in reference
	maxbibnames=9, 			    % Include up to 9 names in citation
	firstinits=true,
]{biblatex}
%%%%%%%%%%%%%%%%%%%%%%%%%%%%%%%%%%%%%%%%%%%%%%%%%%%%%%%%%%%%%%%%%%%%%%%%%%%%%%%%
\addbibresource{../references.bib}
\IfFileExists{bookmark.sty}{\usepackage{bookmark}}{\usepackage{hyperref}}
\IfFileExists{xurl.sty}{\usepackage{xurl}}{} % add URL line breaks if available
\urlstyle{same} % disable monospaced font for URLs
\hypersetup{
  pdftitle={Metropolis-Kimura Evolutionary Dynamics},
  pdfauthor={Manuel Razo-Mejia},
  pdfkeywords={evolutionary dynamics, fitness
landscape, genotype-phenotype density, Metropolis-Hastings, population
genetics},
  colorlinks=true,
  linkcolor={blue},
  filecolor={Maroon},
  citecolor={Blue},
  urlcolor={Blue},
  pdfcreator={LaTeX via pandoc}}

%%%%%%%%%%%%%%%%%%%%%%%%%%%%%%%%%%%%%%%%%%%%%%%%%%%%%%%%%%%%%%%%%%%%%%%%%%%%%%%%
% Add title
\title{Metropolis-Kimura Evolutionary Dynamics}
%%%%%%%%%%%%%%%%%%%%%%%%%%%%%%%%%%%%%%%%%%%%%%%%%%%%%%%%%%%%%%%%%%%%%%%%%%%%%%%%

%%%%%%%%%%%%%%%%%%%%%%%%%%%%%%%%%%%%%%%%%%%%%%%%%%%%%%%%%%%%%%%%%%%%%%%%%%%%%%%%
% Add list of authors

% Loop through authors
\author[1]{Manuel Razo-Mejia}

% Loop through affiliations
\affil[1]{Department of Biology, Stanford University}

% Add corresponding author
%%%%%%%%%%%%%%%%%%%%%%%%%%%%%%%%%%%%%%%%%%%%%%%%%%%%%%%%%%%%%%%%%%%%%%%%%%%%%%%%

%%%%%%%%%%%%%%%%%%%%%%%%%%%%%%%%%%%%%%%%%%%%%%%%%%%%%%%%%%%%%%%%%%%%%%%%%%%%%%%%
% Set affiliations in small font
\renewcommand\Affilfont{\itshape\small}
%%%%%%%%%%%%%%%%%%%%%%%%%%%%%%%%%%%%%%%%%%%%%%%%%%%%%%%%%%%%%%%%%%%%%%%%%%%%%%%%

%%%%%%%%%%%%%%%%%%%%%%%%%%%%%%%%%%%%%%%%%%%%%%%%%%%%%%%%%%%%%%%%%%%%%%%%%%%%%%%%
% Remove date
\date{}
%%%%%%%%%%%%%%%%%%%%%%%%%%%%%%%%%%%%%%%%%%%%%%%%%%%%%%%%%%%%%%%%%%%%%%%%%%%%%%%%

% Add package for line numbering
\usepackage{lineno}

\begin{document}
\linenumbers
\maketitle
\begin{Shaded}
\begin{Highlighting}[]
\CommentTok{\# Import project package}
\ImportTok{import} \BuiltInTok{Antibiotic}

\CommentTok{\# Import basic math libraries}
\ImportTok{import} \BuiltInTok{StatsBase}
\ImportTok{import} \BuiltInTok{LinearAlgebra}
\ImportTok{import} \BuiltInTok{Random}
\ImportTok{import} \BuiltInTok{Distributions}

\CommentTok{\# Load CairoMakie for plotting}
\ImportTok{using} \BuiltInTok{CairoMakie}
\ImportTok{import} \BuiltInTok{ColorSchemes}

\CommentTok{\# Activate backend}
\NormalTok{CairoMakie.}\FunctionTok{activate!}\NormalTok{()}

\CommentTok{\# Set custom plotting style}
\NormalTok{Antibiotic.viz.}\FunctionTok{theme\_makie!}\NormalTok{()}
\end{Highlighting}
\end{Shaded}

\section{Evolutionary dynamics on a fitness and genotype-phenotype
density
landscape}\label{evolutionary-dynamics-on-a-fitness-and-genotype-phenotype-density-landscape}

In this notebook we will explore a simple model of evolution on a
fitness landscape with an explicit genotype-phenotype density function.

The setting is the following:

Let \(\underline{x}(t) = (x_1(t), x_2(t), \cdots, x_N(t))\) be an \(N\)
dimensional vector of phenotypes at time \(t\). The fitness for a given
environment \(E\) is a scalar function \(F_E(\underline{x})\) such that
\begin{equation}\phantomsection\label{eq-fitness-landscape}{
F_E: \mathbb{R}^N \to \mathbb{R},
}\end{equation}

i.e, the coordinates in phenotype space map to a fitness value. As a
particular fitness function we will consider a series of Gaussian
fitness peaks,

\begin{equation}\phantomsection\label{eq-fitness-landscape-gaussian}{
F_E(\underline{x}) = 
\sum_{p=1}^P A_p 
\exp\left(
    -\frac{1}{2}
    (\underline{x} - \underline{\hat{x}}^{(p)})^T 
    \Sigma_p^{-1} 
    (\underline{x} - \underline{\hat{x}}^{(p)})
\right),
}\end{equation}

where \(P\) is the number of peaks, \(A_p\) is the amplitude of the
\(p\)-th peak, \(\underline{\hat{x}}^{(p)}\) is the \(N\)-dimensional
vector of optimal phenotypes for the \(p\)-th peak, and \(\Sigma_p\) is
the \(N \times N\) covariance matrix for the \(p\)-th peak. This
formulation allows for multiple peaks in the fitness landscape, each
potentially having different heights, widths, and covariance structures
between dimensions. The covariance matrix \(\Sigma_p\) captures
potential interactions between different phenotypic dimensions, allowing
for more complex and realistic fitness landscapes.

In the same vein, we define a genotype-phenotype density map
\(GP(\underline{x})\) as a scalar function that maps a phenotype to a
number related to the probability of sampling a given phenotype, i.e.,

\begin{equation}\phantomsection\label{eq-genotype-phenotype-density}{
GP: \mathbb{R}^N \to \mathbb{R},
}\end{equation}

where \(GP(\underline{x})\) is related to the probability of a genotype
mapping to a phenotype \(\underline{x}\). As a particular
genotype-phenotype density function we will consider a set of negative
Gaussian peaks that will serve as ``mutational barriers'' that limit the
range of phenotypes that can be reached. The deeper the peak, the more
difficult it is to reach the phenotype associated with that peak. This
idea is captured by the following genotype-phenotype density function:

\begin{equation}\phantomsection\label{eq-genotype-phenotype-density-gaussian}{
GP(\underline{x}) = 
\sum_{p=1}^P -B_p 
\exp\left(
    -\frac{1}{2}(\underline{x} - \underline{\hat{x}}^{(p)})^T 
    \Sigma_p^{-1} 
    (\underline{x} - \underline{\hat{x}}^{(p)})
\right),
}\end{equation}

where \(B_p\) is the depth of the \(p\)-th peak, and
\(\underline{\hat{x}}^{(p)}\) is the \(N\)-dimensional vector of optimal
phenotypes for the \(p\)-th peak.

Both, the fitness landscape and the genotype-phenotype density function
are built on Gaussian functions. Let's therefore define a
\texttt{struct} to hold the parameters of a single peak.

\begin{Shaded}
\begin{Highlighting}[]
\PreprocessorTok{@doc} \StringTok{raw"""}
\StringTok{    GaussianPeak}

\StringTok{A struct to hold the parameters of a single Gaussian peak.}

\StringTok{\# Fields}
\StringTok{{-} \textasciigrave{}amplitude::AbstractFloat\textasciigrave{}: The amplitude of the peak.}
\StringTok{{-} \textasciigrave{}mean::AbstractVector\textasciigrave{}: The mean of the peak.}
\StringTok{{-} \textasciigrave{}covariance::AbstractMatrix\textasciigrave{}: The covariance matrix of the peak.}
\StringTok{"""}
\KeywordTok{struct}\NormalTok{ GaussianPeak}
\NormalTok{    amplitude}\OperatorTok{::}\DataTypeTok{AbstractFloat}
\NormalTok{    mean}\OperatorTok{::}\DataTypeTok{AbstractVector}
\NormalTok{    covariance}\OperatorTok{::}\DataTypeTok{AbstractMatrix}
\KeywordTok{end}
\end{Highlighting}
\end{Shaded}

Next, let's define a function to evaluate the fitness landscape at a
given phenotype.

\begin{Shaded}
\begin{Highlighting}[]
\PreprocessorTok{@doc} \StringTok{raw"""}
\StringTok{    fitness(peak::GaussianPeak, x::AbstractVecOrMat; min\_value::AbstractFloat=0.0)}
\StringTok{    fitness(peaks::Vector\{GaussianPeak\}, x::AbstractVecOrMat; min\_value::AbstractFloat=0.0)}

\StringTok{Calculate the fitness value for a given phenotype \textasciigrave{}x\textasciigrave{} based on Gaussian peak(s).}

\StringTok{\# Arguments}
\StringTok{{-} \textasciigrave{}peak::GaussianPeak\textasciigrave{}: A single Gaussian peak.}
\StringTok{{-} \textasciigrave{}peaks::Vector\{GaussianPeak\}\textasciigrave{}: A vector of Gaussian peaks.}
\StringTok{{-} \textasciigrave{}x::AbstractVecOrMat\textasciigrave{}: The phenotype(s) for which to calculate the fitness.}
\StringTok{  Can be a vector for a single phenotype or a matrix for multiple phenotypes,}
\StringTok{  where each column corresponds to a phenotype.}
\StringTok{{-} \textasciigrave{}min\_value::AbstractFloat=1.0\textasciigrave{}: The minimum fitness value to be added to the}
\StringTok{  Gaussian contribution.}

\StringTok{\# Returns}
\StringTok{The calculated fitness value(s).}

\StringTok{\# Description}
\StringTok{The first method computes the fitness for a single Gaussian peak, while the}
\StringTok{second method computes the sum of fitness values for multiple Gaussian peaks. In}
\StringTok{both cases, the \textasciigrave{}min\_value\textasciigrave{} is added to the Gaussian contribution to ensure a}
\StringTok{minimum fitness level.}
\StringTok{"""}
\KeywordTok{function} \FunctionTok{fitness}\NormalTok{(}
\NormalTok{    peak}\OperatorTok{::}\DataTypeTok{GaussianPeak}\NormalTok{, x}\OperatorTok{::}\DataTypeTok{AbstractVecOrMat}\NormalTok{; min\_value}\OperatorTok{::}\DataTypeTok{AbstractFloat}\NormalTok{=}\FloatTok{1.0}
\NormalTok{)}
    \CommentTok{\# Calculate the Gaussian peak}
\NormalTok{    gaussian }\OperatorTok{=}\NormalTok{ peak.amplitude }\OperatorTok{*} \FunctionTok{exp}\NormalTok{(}
        \OperatorTok{{-}}\FloatTok{0.5} \OperatorTok{*}
\NormalTok{        (x }\OperatorTok{{-}}\NormalTok{ peak.mean)}\CharTok{\textquotesingle{} *}
        \BuiltInTok{LinearAlgebra}\NormalTok{.}\FunctionTok{inv}\NormalTok{(peak.covariance) }\OperatorTok{*}
\NormalTok{        (x }\OperatorTok{{-}}\NormalTok{ peak.mean)}
\NormalTok{    )}
    \CommentTok{\# Return the fitness by shifting the Gaussian peak by the minimum value}
    \ControlFlowTok{return}\NormalTok{ gaussian }\OperatorTok{+}\NormalTok{ min\_value}
\KeywordTok{end}

\KeywordTok{function} \FunctionTok{fitness}\NormalTok{(}
\NormalTok{    peaks}\OperatorTok{::}\DataTypeTok{Vector\{GaussianPeak\}}\NormalTok{,}
\NormalTok{    x}\OperatorTok{::}\DataTypeTok{AbstractVecOrMat}\NormalTok{;}
\NormalTok{    min\_value}\OperatorTok{::}\DataTypeTok{AbstractFloat}\NormalTok{=}\FloatTok{1.0}
\NormalTok{)}
    \CommentTok{\# Sum the Gaussian contributions from all peaks}
\NormalTok{    total\_gaussian }\OperatorTok{=} \FunctionTok{sum}\NormalTok{(}
\NormalTok{        peak.amplitude }\OperatorTok{*}
        \FunctionTok{exp}\NormalTok{(}
            \OperatorTok{{-}}\FloatTok{0.5} \OperatorTok{*}
\NormalTok{            (x }\OperatorTok{{-}}\NormalTok{ peak.mean)}\CharTok{\textquotesingle{} *}
            \BuiltInTok{LinearAlgebra}\NormalTok{.}\FunctionTok{inv}\NormalTok{(peak.covariance) }\OperatorTok{*}
\NormalTok{            (x }\OperatorTok{{-}}\NormalTok{ peak.mean)}
\NormalTok{        )}
        \ControlFlowTok{for}\NormalTok{ peak }\KeywordTok{in}\NormalTok{ peaks}
\NormalTok{    )}
    \CommentTok{\# Return the fitness by shifting the Gaussian peak by the minimum value}
    \ControlFlowTok{return}\NormalTok{ total\_gaussian }\OperatorTok{+}\NormalTok{ min\_value}
\ControlFlowTok{end}
\end{Highlighting}
\end{Shaded}

Let's test the fitness function with a single Gaussian peak.

\begin{Shaded}
\begin{Highlighting}[]
\CommentTok{\# Define peak parameters}
\NormalTok{amplitude }\OperatorTok{=} \FloatTok{1.0}
\NormalTok{mean }\OperatorTok{=}\NormalTok{ [}\FloatTok{0.0}\NormalTok{, }\FloatTok{0.0}\NormalTok{]}
\NormalTok{covariance }\OperatorTok{=}\NormalTok{ [}\FloatTok{0.5} \FloatTok{0.0}\NormalTok{; }\FloatTok{0.0} \FloatTok{0.5}\NormalTok{]}

\CommentTok{\# Create peak}
\NormalTok{fit\_peak }\OperatorTok{=} \FunctionTok{GaussianPeak}\NormalTok{(amplitude, mean, covariance)}

\CommentTok{\# Define range of phenotypes to evaluate}
\NormalTok{x }\OperatorTok{=} \FunctionTok{range}\NormalTok{(}\OperatorTok{{-}}\FloatTok{2}\NormalTok{, }\FloatTok{2}\NormalTok{, length}\OperatorTok{=}\FloatTok{100}\NormalTok{)}
\NormalTok{y }\OperatorTok{=} \FunctionTok{range}\NormalTok{(}\OperatorTok{{-}}\FloatTok{2}\NormalTok{, }\FloatTok{2}\NormalTok{, length}\OperatorTok{=}\FloatTok{100}\NormalTok{)}

\CommentTok{\# Create meshgrid}
\NormalTok{F }\OperatorTok{=} \FunctionTok{fitness}\NormalTok{.(}\FunctionTok{Ref}\NormalTok{(fit\_peak), [[x, y] for x }\KeywordTok{in}\NormalTok{ x, y }\KeywordTok{in}\NormalTok{ y])}

\CommentTok{\# Initialize figure}
\NormalTok{fig }\OperatorTok{=} \FunctionTok{Figure}\NormalTok{(size}\OperatorTok{=}\NormalTok{(}\FloatTok{700}\NormalTok{, }\FloatTok{300}\NormalTok{))}

\CommentTok{\# Add axis for 3D plot}
\NormalTok{ax3 }\OperatorTok{=} \FunctionTok{Axis3}\NormalTok{(}
\NormalTok{    fig[}\FloatTok{1}\NormalTok{, }\FloatTok{1}\NormalTok{],}
\NormalTok{    xlabel}\OperatorTok{=}\StringTok{"phenotype 1"}\NormalTok{,}
\NormalTok{    ylabel}\OperatorTok{=}\StringTok{"phenotype 2"}\NormalTok{,}
\NormalTok{    zlabel}\OperatorTok{=}\StringTok{"fitness"}\NormalTok{,}
\NormalTok{    yticklabelsvisible}\OperatorTok{=}\ConstantTok{false}\NormalTok{,}
\NormalTok{    xticklabelsvisible}\OperatorTok{=}\ConstantTok{false}\NormalTok{,}
\NormalTok{    zticklabelsvisible}\OperatorTok{=}\ConstantTok{false}\NormalTok{,}
\NormalTok{)}

\CommentTok{\# Add axis for contour plot}
\NormalTok{ax2 }\OperatorTok{=} \FunctionTok{Axis}\NormalTok{(}
\NormalTok{    fig[}\FloatTok{1}\NormalTok{, }\FloatTok{2}\NormalTok{],}
\NormalTok{    xlabel}\OperatorTok{=}\StringTok{"phenotype 1"}\NormalTok{,}
\NormalTok{    ylabel}\OperatorTok{=}\StringTok{"phenotype 2"}\NormalTok{,}
\NormalTok{    aspect}\OperatorTok{=}\FunctionTok{AxisAspect}\NormalTok{(}\FloatTok{1}\NormalTok{),}
\NormalTok{    yticklabelsvisible}\OperatorTok{=}\ConstantTok{false}\NormalTok{,}
\NormalTok{    xticklabelsvisible}\OperatorTok{=}\ConstantTok{false}\NormalTok{,}
\NormalTok{)}

\CommentTok{\# Plot fitness landscape}
\FunctionTok{surface!}\NormalTok{(ax3, x, y, F, color}\OperatorTok{=}\NormalTok{F, colormap}\OperatorTok{=:}\NormalTok{algae)}

\CommentTok{\# Plot a heatmap of the fitness landscape}
\NormalTok{hm }\OperatorTok{=} \FunctionTok{heatmap!}\NormalTok{(ax2, x, y, F, colormap}\OperatorTok{=:}\NormalTok{algae)}

\CommentTok{\# Add a colorbar}
\FunctionTok{Colorbar}\NormalTok{(}
\NormalTok{    fig[}\FloatTok{1}\NormalTok{, }\FloatTok{3}\NormalTok{],}
\NormalTok{    hm,}
\NormalTok{    label}\OperatorTok{=}\StringTok{"fitness"}\NormalTok{,}
\NormalTok{    labelsize}\OperatorTok{=}\FloatTok{13}\NormalTok{,}
\NormalTok{    height}\OperatorTok{=}\FunctionTok{Relative}\NormalTok{(}\FloatTok{1}\NormalTok{),}
\NormalTok{    ticksvisible}\OperatorTok{=}\ConstantTok{false}\NormalTok{,}
\NormalTok{    ticklabelsvisible}\OperatorTok{=}\ConstantTok{false}
\NormalTok{)}
\CommentTok{\# Plot contour plot}
\FunctionTok{contour!}\NormalTok{(ax2, x, y, F, color}\OperatorTok{=:}\NormalTok{white)}

\NormalTok{fig}
\end{Highlighting}
\end{Shaded}

\includegraphics[width=14.58333in,height=6.25in]{metropolis_kimura_evolution_files/figure-pdf/cell-5-output-1.png}

Next, let's define a function to evaluate the genotype-phenotype density
function at a given phenotype.

\begin{Shaded}
\begin{Highlighting}[]
\PreprocessorTok{@doc} \StringTok{raw"""}
\StringTok{    genotype\_phenotype\_density(}
\StringTok{        peak::Union\{GaussianPeak, Vector\{GaussianPeak\}\}, }
\StringTok{        x::AbstractVecOrMat;}
\StringTok{        max\_value::Float64 = 1.0}
\StringTok{    )}

\StringTok{Calculate the genotype{-}phenotype density value for a given phenotype \textasciigrave{}x\textasciigrave{} based}
\StringTok{on Gaussian peak(s).}

\StringTok{\# Arguments}
\StringTok{{-} \textasciigrave{}peak::Union\{GaussianPeak, Vector\{GaussianPeak\}\}\textasciigrave{}: A single Gaussian peak or a}
\StringTok{  vector of Gaussian peaks.}
\StringTok{{-} \textasciigrave{}x::AbstractVecOrMat\textasciigrave{}: The phenotype(s) for which to calculate the genotype{-}}
\StringTok{  phenotype density. Can be a vector for a single phenotype or a matrix for}
\StringTok{  multiple phenotypes, where each column corresponds to a phenotype.}
\StringTok{{-} \textasciigrave{}max\_value::Float64\textasciigrave{}: Optional. The maximum value of the genotype{-}phenotype}
\StringTok{  density. Default is 1.0.}
\StringTok{  1.0.}

\StringTok{\# Returns}
\StringTok{The calculated genotype{-}phenotype density value(s).}
\StringTok{"""}
\KeywordTok{function} \FunctionTok{genotype\_phenotype\_density}\NormalTok{(}
\NormalTok{    peak}\OperatorTok{::}\DataTypeTok{GaussianPeak}\NormalTok{, x}\OperatorTok{::}\DataTypeTok{AbstractVecOrMat}\NormalTok{; max\_value}\OperatorTok{::}\DataTypeTok{AbstractFloat}\NormalTok{=}\FloatTok{1.0}
\NormalTok{)}
    \CommentTok{\# Calculate the Gaussian peak}
\NormalTok{    gaussian }\OperatorTok{=}\NormalTok{ peak.amplitude }\OperatorTok{*} \FunctionTok{exp}\NormalTok{(}
        \OperatorTok{{-}}\FloatTok{0.5} \OperatorTok{*}
\NormalTok{        (x }\OperatorTok{{-}}\NormalTok{ peak.mean)}\CharTok{\textquotesingle{} *}
        \BuiltInTok{LinearAlgebra}\NormalTok{.}\FunctionTok{inv}\NormalTok{(peak.covariance) }\OperatorTok{*}
\NormalTok{        (x }\OperatorTok{{-}}\NormalTok{ peak.mean)}
\NormalTok{    )}
    \CommentTok{\# Return the genotype{-}phenotype density by inverting the Gaussian peak}
    \CommentTok{\# shifted by the maximum value}
    \ControlFlowTok{return} \FunctionTok{{-}}\NormalTok{(gaussian }\OperatorTok{{-}}\NormalTok{ max\_value)}
\KeywordTok{end}

\KeywordTok{function} \FunctionTok{genotype\_phenotype\_density}\NormalTok{(}
\NormalTok{    peaks}\OperatorTok{::}\DataTypeTok{Vector\{GaussianPeak\}}\NormalTok{,}
\NormalTok{    x}\OperatorTok{::}\DataTypeTok{AbstractVecOrMat}\NormalTok{;}
\NormalTok{    max\_value}\OperatorTok{::}\DataTypeTok{AbstractFloat}\NormalTok{=}\FloatTok{1.0}
\NormalTok{)}
    \CommentTok{\# Sum the Gaussian contributions from all peaks}
\NormalTok{    total\_gaussian }\OperatorTok{=} \FunctionTok{sum}\NormalTok{(}
\NormalTok{        peak.amplitude }\OperatorTok{*}
        \FunctionTok{exp}\NormalTok{(}
            \OperatorTok{{-}}\FloatTok{0.5} \OperatorTok{*}
\NormalTok{            (x }\OperatorTok{{-}}\NormalTok{ peak.mean)}\CharTok{\textquotesingle{} *}
            \BuiltInTok{LinearAlgebra}\NormalTok{.}\FunctionTok{inv}\NormalTok{(peak.covariance) }\OperatorTok{*}
\NormalTok{            (x }\OperatorTok{{-}}\NormalTok{ peak.mean)}
\NormalTok{        )}
        \ControlFlowTok{for}\NormalTok{ peak }\KeywordTok{in}\NormalTok{ peaks}
\NormalTok{    )}
    \CommentTok{\# Return the genotype{-}phenotype density by inverting the Gaussian peak}
    \CommentTok{\# shifted by the maximum value}
    \ControlFlowTok{return} \FunctionTok{{-}}\NormalTok{(total\_gaussian }\OperatorTok{{-}}\NormalTok{ max\_value)}
\ControlFlowTok{end}
\end{Highlighting}
\end{Shaded}

Let's test the genotype-phenotype density function with four Gaussian
peaks in the corners of a square in phenotype space.

\begin{Shaded}
\begin{Highlighting}[]
\CommentTok{\# Define peak parameters}
\NormalTok{amplitude }\OperatorTok{=} \FloatTok{0.5}
\NormalTok{means }\OperatorTok{=}\NormalTok{ [}
\NormalTok{    [}\OperatorTok{{-}}\FloatTok{1.5}\NormalTok{, }\OperatorTok{{-}}\FloatTok{1.5}\NormalTok{],}
\NormalTok{    [}\FloatTok{1.5}\NormalTok{, }\OperatorTok{{-}}\FloatTok{1.5}\NormalTok{],}
\NormalTok{    [}\FloatTok{1.5}\NormalTok{, }\FloatTok{1.5}\NormalTok{],}
\NormalTok{    [}\OperatorTok{{-}}\FloatTok{1.5}\NormalTok{, }\FloatTok{1.5}\NormalTok{],}
\NormalTok{]}
\NormalTok{covariance }\OperatorTok{=}\NormalTok{ [}\FloatTok{0.45} \FloatTok{0.0}\NormalTok{; }\FloatTok{0.0} \FloatTok{0.45}\NormalTok{]}

\CommentTok{\# Create peak}
\NormalTok{mut\_peaks }\OperatorTok{=} \FunctionTok{GaussianPeak}\NormalTok{.(}\FunctionTok{Ref}\NormalTok{(amplitude), means, }\FunctionTok{Ref}\NormalTok{(covariance))}

\CommentTok{\# Define range of phenotypes to evaluate}
\NormalTok{x }\OperatorTok{=} \FunctionTok{range}\NormalTok{(}\OperatorTok{{-}}\FloatTok{3}\NormalTok{, }\FloatTok{3}\NormalTok{, length}\OperatorTok{=}\FloatTok{100}\NormalTok{)}
\NormalTok{y }\OperatorTok{=} \FunctionTok{range}\NormalTok{(}\OperatorTok{{-}}\FloatTok{3}\NormalTok{, }\FloatTok{3}\NormalTok{, length}\OperatorTok{=}\FloatTok{100}\NormalTok{)}

\CommentTok{\# Create meshgrid}
\NormalTok{M }\OperatorTok{=} \FunctionTok{genotype\_phenotype\_density}\NormalTok{.(}\FunctionTok{Ref}\NormalTok{(mut\_peaks), [[x, y] for x }\KeywordTok{in}\NormalTok{ x, y }\KeywordTok{in}\NormalTok{ y])}

\CommentTok{\# Initialize figure}
\NormalTok{fig }\OperatorTok{=} \FunctionTok{Figure}\NormalTok{(size}\OperatorTok{=}\NormalTok{(}\FloatTok{700}\NormalTok{, }\FloatTok{300}\NormalTok{))}

\CommentTok{\# Add axis for 3D plot}
\NormalTok{ax3 }\OperatorTok{=} \FunctionTok{Axis3}\NormalTok{(}
\NormalTok{    fig[}\FloatTok{1}\NormalTok{, }\FloatTok{1}\NormalTok{],}
\NormalTok{    xlabel}\OperatorTok{=}\StringTok{"phenotype 1"}\NormalTok{,}
\NormalTok{    ylabel}\OperatorTok{=}\StringTok{"phenotype 2"}\NormalTok{,}
\NormalTok{    zlabel}\OperatorTok{=}\StringTok{"∝ probability of }\SpecialCharTok{\textbackslash{}n}\StringTok{phenotype"}\NormalTok{,}
\NormalTok{    yticklabelsvisible}\OperatorTok{=}\ConstantTok{false}\NormalTok{,}
\NormalTok{    xticklabelsvisible}\OperatorTok{=}\ConstantTok{false}\NormalTok{,}
\NormalTok{    zticklabelsvisible}\OperatorTok{=}\ConstantTok{false}\NormalTok{,}
\NormalTok{)}

\CommentTok{\# Add axis for contour plot}
\NormalTok{ax2 }\OperatorTok{=} \FunctionTok{Axis}\NormalTok{(}
\NormalTok{    fig[}\FloatTok{1}\NormalTok{, }\FloatTok{2}\NormalTok{],}
\NormalTok{    xlabel}\OperatorTok{=}\StringTok{"phenotype 1"}\NormalTok{,}
\NormalTok{    ylabel}\OperatorTok{=}\StringTok{"phenotype 2"}\NormalTok{,}
\NormalTok{    aspect}\OperatorTok{=}\FunctionTok{AxisAspect}\NormalTok{(}\FloatTok{1}\NormalTok{),}
\NormalTok{    yticklabelsvisible}\OperatorTok{=}\ConstantTok{false}\NormalTok{,}
\NormalTok{    xticklabelsvisible}\OperatorTok{=}\ConstantTok{false}\NormalTok{,}
\NormalTok{)}

\CommentTok{\# Plot genotype{-}phenotype density}
\FunctionTok{surface!}\NormalTok{(ax3, x, y, M, color}\OperatorTok{=}\NormalTok{M, colormap}\OperatorTok{=}\FunctionTok{Reverse}\NormalTok{(ColorSchemes.Purples\_9))}

\CommentTok{\# Plot a heatmap of the genotype{-}phenotype density}
\NormalTok{hm }\OperatorTok{=} \FunctionTok{heatmap!}\NormalTok{(ax2, x, y, M, colormap}\OperatorTok{=}\FunctionTok{Reverse}\NormalTok{(ColorSchemes.Purples\_9))}

\CommentTok{\# Add a colorbar}
\FunctionTok{Colorbar}\NormalTok{(}
\NormalTok{    fig[}\FloatTok{1}\NormalTok{, }\FloatTok{3}\NormalTok{],}
\NormalTok{    hm,}
\NormalTok{    label}\OperatorTok{=}\StringTok{"∝ probability of phenotype"}\NormalTok{,}
\NormalTok{    labelsize}\OperatorTok{=}\FloatTok{13}\NormalTok{,}
\NormalTok{    height}\OperatorTok{=}\FunctionTok{Relative}\NormalTok{(}\FloatTok{1}\NormalTok{),}
\NormalTok{    ticksvisible}\OperatorTok{=}\ConstantTok{false}\NormalTok{,}
\NormalTok{    ticklabelsvisible}\OperatorTok{=}\ConstantTok{false}
\NormalTok{)}

\CommentTok{\# Plot contour plot}
\FunctionTok{contour!}\NormalTok{(ax2, x, y, M, color}\OperatorTok{=:}\NormalTok{white)}

\NormalTok{fig}
\end{Highlighting}
\end{Shaded}

\includegraphics[width=14.58333in,height=6.25in]{metropolis_kimura_evolution_files/figure-pdf/cell-7-output-1.png}

\subsection{Evolutionary dynamics}\label{evolutionary-dynamics}

\subsubsection{Metropolis-Kimura Algorithm in Evolutionary
Dynamics}\label{metropolis-kimura-algorithm-in-evolutionary-dynamics}

We now introduce a biologically motivated algorithm that combines
elements of the Metropolis-Hastings algorithm with Motoo Kimura's
population genetics theory. This approach implements a two-step process
that separately models:

\begin{enumerate}
\def\labelenumi{\arabic{enumi}.}
\tightlist
\item
  The probability of mutation occurring (mutation accessibility) as
  governed by the genotype-phenotype density.
\item
  The probability of fixation within a population (selection) determined
  by the fitness effect of such mutation and the effective population
  size.
\end{enumerate}

\subsubsection{Biological Motivation}\label{biological-motivation}

In natural populations, evolution proceeds through two distinct
processes:

\begin{enumerate}
\def\labelenumi{\arabic{enumi}.}
\item
  \textbf{Mutation}: New variants arise through random genetic changes.
  The probability of specific mutations depends on molecular mechanisms
  and constraints of the genotype-phenotype map.
\item
  \textbf{Fixation}: Once a mutation occurs, it must spread through the
  population to become fixed. The probability of fixation depends on the
  selection coefficient (fitness advantage) and the effective population
  size.
\end{enumerate}

\subsubsection{Mathematical Framework}\label{mathematical-framework}

\paragraph{Mutation Probability}\label{mutation-probability}

The probability of a mutation from phenotype \(\underline{x}\) to
\(\underline{x}'\) is modeled using a Metropolis-like criterion based on
the genotype-phenotype density:

\begin{equation}\phantomsection\label{eq-mutation-probability}{
\pi_{\text{mut}}(\underline{x} \to \underline{x}') = 
\min\left(1, \frac{GP(\underline{x}')}{GP(\underline{x})}\right)^{\beta}
}\end{equation}

Here, \(GP(\underline{x})\) represents the genotype-phenotype density at
phenotype \(\underline{x}\), and \(\beta\) is a parameter controlling
the strength of mutational constraints. The higher \(\beta\) is, the
less likely a mutation going downhill in genotype-phenotype density is
to be accepted.

\paragraph{Kimura's Fixation
Probability}\label{kimuras-fixation-probability}

Once a mutation occurs, its probability of fixation in a population of
effective size \(N\) is given by Kimura's formula

\begin{equation}\phantomsection\label{eq-kimura-fixation-probability}{
\pi_{\text{fix}}(\underline{x} \to \underline{x}') = 
\frac{1 - e^{-2s}}{1 - e^{-2Ns}}
}\end{equation}

where \(s\) is the selection coefficient. We compute this selection
coefficient as the difference in fitness between the new and old
phenotype divided by the fitness of the old phenotype, i.e.,

\begin{equation}\phantomsection\label{eq-selection-coefficient}{
s = \frac{F_E(\underline{x}') - F_E(\underline{x})}{F_E(\underline{x})}
}\end{equation}

This equation captures a fundamental result from population genetics:
beneficial mutations (\(s > 0\)) have a higher probability of fixation,
with the probability approaching \(2s\) for small positive selection
coefficients and approaching 1 for large positive selection
coefficients. Deleterious mutations (\(s < 0\)) have an exponentially
decreasing probability of fixation as population size increases.

\paragraph{Overall Acceptance Probability of
Mutation}\label{overall-acceptance-probability-of-mutation}

The overall acceptance probability---defined as the probability of a
proposed step \(x \rightarrow x'\)---is the product of the mutation
probability and the fixation probability

\begin{equation}\phantomsection\label{eq-overall-acceptance-probability}{
\pi_{\text{accept}}(\underline{x} \to \underline{x}') = 
\pi_{\text{mut}}(\underline{x} \to \underline{x}') \cdot 
\pi_{\text{fix}}(\underline{x} \to \underline{x}')
}\end{equation}

This two-step process better reflects the biological reality of
evolution, where both mutational accessibility and selection contribute
to evolutionary trajectories.

\subsubsection{Numerical Implementation}\label{numerical-implementation}

Let's now implement the Metropolis-Kimura algorithm in code.

\begin{Shaded}
\begin{Highlighting}[]
\PreprocessorTok{@doc} \StringTok{raw"""}
\StringTok{    evo\_metropolis\_kimura(x0, fitness\_peaks, gen\_peaks, N, β, µ, n\_steps)}

\StringTok{Perform evolutionary algorithm to simulate phenotypic evolution using a}
\StringTok{Metropolis{-}Hastings{-}like mutation probability and Kimura\textquotesingle{}s fixation probability.}

\StringTok{\# Arguments}
\StringTok{{-} \textasciigrave{}x0::AbstractVector\textasciigrave{}: Initial phenotype vector.}
\StringTok{{-} \textasciigrave{}fitness\_peaks::Union\{GaussianPeak, Vector\{GaussianPeak\}\}\textasciigrave{}: Fitness landscape}
\StringTok{  defined by one or more Gaussian peaks.}
\StringTok{{-} \textasciigrave{}gen\_peaks::Union\{GaussianPeak, Vector\{GaussianPeak\}\}\textasciigrave{}: Genotype{-}phenotype}
\StringTok{  density landscape defined by one or more Gaussian peaks.}
\StringTok{{-} \textasciigrave{}N::Real\textasciigrave{}: Effective population size.}
\StringTok{{-} \textasciigrave{}β::Real\textasciigrave{}: Parameter controlling the strength of mutational constraints.}
\StringTok{{-} \textasciigrave{}µ::Real\textasciigrave{}: Mutation step size standard deviation.}
\StringTok{{-} \textasciigrave{}n\_steps::Int\textasciigrave{}: Number of steps to simulate.}

\StringTok{\# Returns}
\StringTok{{-} \textasciigrave{}Matrix\{Float64\}\textasciigrave{}: Matrix of phenotypes, where each column represents a step}
\StringTok{  in the simulation.}

\StringTok{\# Description}
\StringTok{This function implements a two{-}step evolutionary algorithm to simulate}
\StringTok{phenotypic evolution in a landscape defined by fitness and genotype{-}phenotype}
\StringTok{density. It uses the following steps:}

\StringTok{1. Initialize the phenotype trajectory with the given starting point.}
\StringTok{2. For each step: }
\StringTok{   a. Propose a new phenotype by adding Gaussian noise. }
\StringTok{   b. Calculate the fitness and genotype{-}phenotype density values for the new }
\StringTok{      phenotype.}
\StringTok{   c. Compute the mutation probability P\_mut = min(1, (GP(x\_new)/GP(x))\^{}β)}
\StringTok{   d. Compute the selection coefficient s = (F\_new {-} F)/F}
\StringTok{   e. Calculate fixation probability using Kimura\textquotesingle{}s equation:}
\StringTok{      P\_fix = (1 {-} exp({-}2s))/(1 {-} exp({-}2Ns))}
\StringTok{   f. Accept or reject based on P\_accept = P\_mut * P\_fix}
\StringTok{3. Return the complete phenotype trajectory.}
\StringTok{"""}
\KeywordTok{function} \FunctionTok{evo\_metropolis\_kimura}\NormalTok{(}
\NormalTok{    x0}\OperatorTok{::}\DataTypeTok{AbstractVector}\NormalTok{,}
\NormalTok{    fitness\_peaks}\OperatorTok{::}\DataTypeTok{Union\{GaussianPeak,Vector\{GaussianPeak\}\}}\NormalTok{,}
\NormalTok{    gen\_peaks}\OperatorTok{::}\DataTypeTok{Union\{GaussianPeak,Vector\{GaussianPeak\}\}}\NormalTok{,}
\NormalTok{    N}\OperatorTok{::}\DataTypeTok{Real}\NormalTok{,}
\NormalTok{    β}\OperatorTok{::}\DataTypeTok{Real}\NormalTok{,}
\NormalTok{    µ}\OperatorTok{::}\DataTypeTok{Real}\NormalTok{,}
\NormalTok{    n\_steps}\OperatorTok{::}\DataTypeTok{Int}\NormalTok{,}
\NormalTok{)}
    \CommentTok{\# Initialize array to hold phenotypes}
\NormalTok{    x }\OperatorTok{=} \FunctionTok{Matrix}\DataTypeTok{\{Float64\}}\NormalTok{(}\ConstantTok{undef}\NormalTok{, }\FunctionTok{length}\NormalTok{(x0), n\_steps }\OperatorTok{+} \FloatTok{1}\NormalTok{)}
    \CommentTok{\# Set initial phenotype}
\NormalTok{    x[}\OperatorTok{:}\NormalTok{, }\FloatTok{1}\NormalTok{] }\OperatorTok{=}\NormalTok{ x0}

    \CommentTok{\# Compute fitness and genotype{-}phenotype density at initial phenotype}
\NormalTok{    fitness\_val }\OperatorTok{=} \FunctionTok{fitness}\NormalTok{(fitness\_peaks, x0)}
\NormalTok{    gen\_val }\OperatorTok{=} \FunctionTok{genotype\_phenotype\_density}\NormalTok{(gen\_peaks, x0)}

    \CommentTok{\# Loop over steps}
    \ControlFlowTok{for}\NormalTok{ t }\KeywordTok{in} \FloatTok{1}\OperatorTok{:}\NormalTok{n\_steps}
        \CommentTok{\# Propose new phenotype}
\NormalTok{        x\_new }\OperatorTok{=}\NormalTok{ x[}\OperatorTok{:}\NormalTok{, t] }\OperatorTok{+}\NormalTok{ µ }\OperatorTok{*} \FunctionTok{randn}\NormalTok{(}\FunctionTok{length}\NormalTok{(x0))}

        \CommentTok{\# Calculate fitness and genotype{-}phenotype density}
\NormalTok{        fitness\_val\_new }\OperatorTok{=} \FunctionTok{fitness}\NormalTok{(fitness\_peaks, x\_new)}
\NormalTok{        gen\_val\_new }\OperatorTok{=} \FunctionTok{genotype\_phenotype\_density}\NormalTok{(gen\_peaks, x\_new)}

        \CommentTok{\# Compute selection coefficient}
\NormalTok{        s }\OperatorTok{=}\NormalTok{ (fitness\_val\_new }\OperatorTok{{-}}\NormalTok{ fitness\_val) }\OperatorTok{/}
\NormalTok{            (fitness\_val }\OperatorTok{+} \FunctionTok{eps}\NormalTok{(}\FunctionTok{eltype}\NormalTok{(fitness\_val)))}

        \CommentTok{\# Compute mutation probability}
\NormalTok{        P\_mut }\OperatorTok{=} \FunctionTok{min}\NormalTok{(}\FloatTok{1}\NormalTok{, (gen\_val\_new }\OperatorTok{/}\NormalTok{ gen\_val)}\OperatorTok{\^{}}\NormalTok{β)}

        \CommentTok{\# Compute fixation probability using Kimura\textquotesingle{}s equation}
        \CommentTok{\# Use approximation when s is very small for numerical stability}
        \ControlFlowTok{if} \FunctionTok{abs}\NormalTok{(s) }\OperatorTok{\textless{}} \FloatTok{1e{-}10}
\NormalTok{            P\_fix }\OperatorTok{=} \FloatTok{1} \OperatorTok{/}\NormalTok{ N}
        \ControlFlowTok{elseif} \FunctionTok{abs}\NormalTok{(s) }\OperatorTok{\textless{}} \FloatTok{0.1}
            \CommentTok{\# For small s, use approximation}
\NormalTok{            P\_fix }\OperatorTok{=}\NormalTok{ (}\FloatTok{1} \OperatorTok{{-}} \FunctionTok{exp}\NormalTok{(}\OperatorTok{{-}}\FloatTok{2} \OperatorTok{*}\NormalTok{ s)) }\OperatorTok{/}\NormalTok{ (}\FloatTok{1} \OperatorTok{{-}} \FunctionTok{exp}\NormalTok{(}\OperatorTok{{-}}\FloatTok{2} \OperatorTok{*}\NormalTok{ N }\OperatorTok{*}\NormalTok{ s))}
            \CommentTok{\# Check for potential numerical issues}
            \ControlFlowTok{if}\NormalTok{ !}\FunctionTok{isfinite}\NormalTok{(P\_fix) }\OperatorTok{||}\NormalTok{ P\_fix }\OperatorTok{\textless{}} \FloatTok{0}
\NormalTok{                P\_fix }\OperatorTok{=}\NormalTok{ s }\OperatorTok{\textgreater{}} \FloatTok{0}\NormalTok{ ? }\FloatTok{2} \OperatorTok{*}\NormalTok{ s }\OperatorTok{:} \FloatTok{1} \OperatorTok{/}\NormalTok{ N}
            \ControlFlowTok{end}
        \ControlFlowTok{else}
            \CommentTok{\# Otherwise use full Kimura equation}
\NormalTok{            P\_fix }\OperatorTok{=}\NormalTok{ (}\FloatTok{1} \OperatorTok{{-}} \FunctionTok{exp}\NormalTok{(}\OperatorTok{{-}}\FloatTok{2} \OperatorTok{*}\NormalTok{ s)) }\OperatorTok{/}\NormalTok{ (}\FloatTok{1} \OperatorTok{{-}} \FunctionTok{exp}\NormalTok{(}\OperatorTok{{-}}\FloatTok{2} \OperatorTok{*}\NormalTok{ N }\OperatorTok{*}\NormalTok{ s))}
        \ControlFlowTok{end}

        \CommentTok{\# Compute acceptance probability}
\NormalTok{        P\_accept }\OperatorTok{=}\NormalTok{ P\_mut }\OperatorTok{*}\NormalTok{ P\_fix}

        \CommentTok{\# Accept or reject proposal}
        \ControlFlowTok{if} \FunctionTok{rand}\NormalTok{() }\OperatorTok{\textless{}}\NormalTok{ P\_accept}
            \CommentTok{\# Accept proposal}
\NormalTok{            x[}\OperatorTok{:}\NormalTok{, t}\OperatorTok{+}\FloatTok{1}\NormalTok{] }\OperatorTok{=}\NormalTok{ x\_new}
            \CommentTok{\# Update fitness and genotype{-}phenotype density}
\NormalTok{            fitness\_val }\OperatorTok{=}\NormalTok{ fitness\_val\_new}
\NormalTok{            gen\_val }\OperatorTok{=}\NormalTok{ gen\_val\_new}
        \ControlFlowTok{else}
            \CommentTok{\# Reject proposal}
\NormalTok{            x[}\OperatorTok{:}\NormalTok{, t}\OperatorTok{+}\FloatTok{1}\NormalTok{] }\OperatorTok{=}\NormalTok{ x[}\OperatorTok{:}\NormalTok{, t]}
        \ControlFlowTok{end}
    \ControlFlowTok{end} \CommentTok{\# for}
    \ControlFlowTok{return}\NormalTok{ x}
\KeywordTok{end} \CommentTok{\# function}
\end{Highlighting}
\end{Shaded}

Let's test this function with the previously defined fitness and
genotype-phenotype density functions.

\begin{Shaded}
\begin{Highlighting}[]
\BuiltInTok{Random}\NormalTok{.}\FunctionTok{seed!}\NormalTok{(}\FloatTok{42}\NormalTok{)}

\CommentTok{\# Define peak parameters}
\NormalTok{amplitude }\OperatorTok{=} \FloatTok{5.0}
\NormalTok{mean }\OperatorTok{=}\NormalTok{ [}\FloatTok{0.0}\NormalTok{, }\FloatTok{0.0}\NormalTok{]}
\NormalTok{covariance }\OperatorTok{=}\NormalTok{ [}\FloatTok{3.0} \FloatTok{0.0}\NormalTok{; }\FloatTok{0.0} \FloatTok{3.0}\NormalTok{]}
\CommentTok{\# Create peak}
\NormalTok{fit\_peak }\OperatorTok{=} \FunctionTok{GaussianPeak}\NormalTok{(amplitude, mean, covariance)}
\CommentTok{\# Define peak parameters}
\NormalTok{amplitude }\OperatorTok{=} \FloatTok{1.0}
\NormalTok{means }\OperatorTok{=}\NormalTok{ [}
\NormalTok{    [}\OperatorTok{{-}}\FloatTok{1.5}\NormalTok{, }\OperatorTok{{-}}\FloatTok{1.5}\NormalTok{],}
\NormalTok{    [}\FloatTok{1.5}\NormalTok{, }\OperatorTok{{-}}\FloatTok{1.5}\NormalTok{],}
\NormalTok{    [}\FloatTok{1.5}\NormalTok{, }\FloatTok{1.5}\NormalTok{],}
\NormalTok{    [}\OperatorTok{{-}}\FloatTok{1.5}\NormalTok{, }\FloatTok{1.5}\NormalTok{],}
\NormalTok{]}
\NormalTok{covariance }\OperatorTok{=}\NormalTok{ [}\FloatTok{0.45} \FloatTok{0.0}\NormalTok{; }\FloatTok{0.0} \FloatTok{0.45}\NormalTok{]}

\CommentTok{\# Create peak}
\NormalTok{mut\_peaks }\OperatorTok{=} \FunctionTok{GaussianPeak}\NormalTok{.(}\FunctionTok{Ref}\NormalTok{(amplitude), means, }\FunctionTok{Ref}\NormalTok{(covariance))}

\CommentTok{\# Set initial phenotype}
\NormalTok{x0 }\OperatorTok{=}\NormalTok{ [}\OperatorTok{{-}}\FloatTok{2.5}\NormalTok{, }\OperatorTok{{-}}\FloatTok{2.5}\NormalTok{]}

\CommentTok{\# Set parameters}
\NormalTok{N }\OperatorTok{=} \FloatTok{10} \CommentTok{\# Effective population size }
\NormalTok{β }\OperatorTok{=} \FloatTok{10.0}    \CommentTok{\# Mutational constraint parameter}
\NormalTok{µ }\OperatorTok{=} \FloatTok{0.1}    \CommentTok{\# Mutation step size}
\NormalTok{n\_steps }\OperatorTok{=} \FloatTok{3000}

\CommentTok{\# Run Metropolis{-}Kimura algorithm}
\NormalTok{x\_traj }\OperatorTok{=} \FunctionTok{evo\_metropolis\_kimura}\NormalTok{(x0, fit\_peak, mut\_peaks, N, β, µ, n\_steps)}
\end{Highlighting}
\end{Shaded}

Let's plot the trajectory of the phenotypes in phenotype space.

\begin{Shaded}
\begin{Highlighting}[]
\BuiltInTok{Random}\NormalTok{.}\FunctionTok{seed!}\NormalTok{(}\FloatTok{42}\NormalTok{)}

\CommentTok{\# Define range of phenotypes to evaluate}
\NormalTok{x }\OperatorTok{=} \FunctionTok{range}\NormalTok{(}\OperatorTok{{-}}\FloatTok{4}\NormalTok{, }\FloatTok{4}\NormalTok{, length}\OperatorTok{=}\FloatTok{100}\NormalTok{)}
\NormalTok{y }\OperatorTok{=} \FunctionTok{range}\NormalTok{(}\OperatorTok{{-}}\FloatTok{4}\NormalTok{, }\FloatTok{4}\NormalTok{, length}\OperatorTok{=}\FloatTok{100}\NormalTok{)}

\CommentTok{\# Create meshgrid}
\NormalTok{F }\OperatorTok{=} \FunctionTok{fitness}\NormalTok{.(}\FunctionTok{Ref}\NormalTok{(fit\_peak), [[x, y] for x }\KeywordTok{in}\NormalTok{ x, y }\KeywordTok{in}\NormalTok{ y])}
\NormalTok{M }\OperatorTok{=} \FunctionTok{genotype\_phenotype\_density}\NormalTok{.(}\FunctionTok{Ref}\NormalTok{(mut\_peaks), [[x, y] for x }\KeywordTok{in}\NormalTok{ x, y }\KeywordTok{in}\NormalTok{ y])}

\CommentTok{\# Initialize figure}
\NormalTok{fig }\OperatorTok{=} \FunctionTok{Figure}\NormalTok{(size}\OperatorTok{=}\NormalTok{(}\FloatTok{600}\NormalTok{, }\FloatTok{300}\NormalTok{))}

\CommentTok{\# Add axis for trajectory in fitness landscape}
\NormalTok{ax1 }\OperatorTok{=} \FunctionTok{Axis}\NormalTok{(}
\NormalTok{    fig[}\FloatTok{1}\NormalTok{, }\FloatTok{1}\NormalTok{],}
\NormalTok{    xlabel}\OperatorTok{=}\StringTok{"phenotype 1"}\NormalTok{,}
\NormalTok{    ylabel}\OperatorTok{=}\StringTok{"phenotype 2"}\NormalTok{,}
\NormalTok{    aspect}\OperatorTok{=}\FunctionTok{AxisAspect}\NormalTok{(}\FloatTok{1}\NormalTok{),}
\NormalTok{    title}\OperatorTok{=}\StringTok{"Fitness landscape"}\NormalTok{,}
\NormalTok{    yticklabelsvisible}\OperatorTok{=}\ConstantTok{false}\NormalTok{,}
\NormalTok{    xticklabelsvisible}\OperatorTok{=}\ConstantTok{false}\NormalTok{,}
\NormalTok{)}
\CommentTok{\# Add axis for trajectory in genotype{-}phenotype density}
\NormalTok{ax2 }\OperatorTok{=} \FunctionTok{Axis}\NormalTok{(}
\NormalTok{    fig[}\FloatTok{1}\NormalTok{, }\FloatTok{2}\NormalTok{],}
\NormalTok{    xlabel}\OperatorTok{=}\StringTok{"phenotype 1"}\NormalTok{,}
\NormalTok{    ylabel}\OperatorTok{=}\StringTok{"phenotype 2"}\NormalTok{,}
\NormalTok{    aspect}\OperatorTok{=}\FunctionTok{AxisAspect}\NormalTok{(}\FloatTok{1}\NormalTok{),}
\NormalTok{    title}\OperatorTok{=}\StringTok{"Genotype{-}phenotype}\SpecialCharTok{\textbackslash{}n}\StringTok{Density"}\NormalTok{,}
\NormalTok{    yticklabelsvisible}\OperatorTok{=}\ConstantTok{false}\NormalTok{,}
\NormalTok{    xticklabelsvisible}\OperatorTok{=}\ConstantTok{false}\NormalTok{,}
\NormalTok{)}

\CommentTok{\# Plot a heatmap of the fitness landscape}
\FunctionTok{heatmap!}\NormalTok{(ax1, x, y, F, colormap}\OperatorTok{=:}\NormalTok{algae)}
\CommentTok{\# Plot heatmap of genotype{-}phenotype density}
\FunctionTok{heatmap!}\NormalTok{(ax2, x, y, M, colormap}\OperatorTok{=}\FunctionTok{Reverse}\NormalTok{(ColorSchemes.Purples\_9))}

\CommentTok{\# Plot contour plot}
\FunctionTok{contour!}\NormalTok{(ax1, x, y, F, color}\OperatorTok{=:}\NormalTok{white)}
\FunctionTok{contour!}\NormalTok{(ax2, x, y, M, color}\OperatorTok{=:}\NormalTok{white)}

\CommentTok{\# Plot trajectory}
\FunctionTok{scatterlines!}\NormalTok{.(}
\NormalTok{    [ax1, ax2],}
    \FunctionTok{Ref}\NormalTok{(x\_traj[}\FloatTok{1}\NormalTok{, }\OperatorTok{:}\NormalTok{]),}
    \FunctionTok{Ref}\NormalTok{(x\_traj[}\FloatTok{2}\NormalTok{, }\OperatorTok{:}\NormalTok{]),}
\NormalTok{    color}\OperatorTok{=}\NormalTok{ColorSchemes.seaborn\_colorblind[}\FloatTok{4}\NormalTok{],}
\NormalTok{    markersize}\OperatorTok{=}\FloatTok{3}
\NormalTok{)}

\NormalTok{fig}
\end{Highlighting}
\end{Shaded}

\includegraphics[width=12.5in,height=6.25in]{metropolis_kimura_evolution_files/figure-pdf/cell-10-output-1.png}

We can see that the trajectory climbs the fitness peak, but it does so
by avoiding regions of low genotype-phenotype density.

Let's plot the fitness and the genotype-phenotype density for the
phenotypes in the trajectory over time.

\begin{Shaded}
\begin{Highlighting}[]
\CommentTok{\# Initialize arrays to hold fitness and genotype{-}phenotype density}
\NormalTok{F\_traj }\OperatorTok{=} \FunctionTok{Vector}\DataTypeTok{\{Float64\}}\NormalTok{(}\ConstantTok{undef}\NormalTok{, n\_steps }\OperatorTok{+} \FloatTok{1}\NormalTok{)}
\NormalTok{M\_traj }\OperatorTok{=} \FunctionTok{Vector}\DataTypeTok{\{Float64\}}\NormalTok{(}\ConstantTok{undef}\NormalTok{, n\_steps }\OperatorTok{+} \FloatTok{1}\NormalTok{)}

\CommentTok{\# Loop over steps}
\ControlFlowTok{for}\NormalTok{ (i, x) }\KeywordTok{in} \FunctionTok{enumerate}\NormalTok{(}\FunctionTok{eachcol}\NormalTok{(x\_traj))}
\NormalTok{    F\_traj[i] }\OperatorTok{=} \FunctionTok{fitness}\NormalTok{(fit\_peak, x)}
\NormalTok{    M\_traj[i] }\OperatorTok{=} \FunctionTok{genotype\_phenotype\_density}\NormalTok{(mut\_peaks, x)}
\ControlFlowTok{end}

\CommentTok{\# Initialize figure}
\NormalTok{fig }\OperatorTok{=} \FunctionTok{Figure}\NormalTok{(size}\OperatorTok{=}\NormalTok{(}\FloatTok{600}\NormalTok{, }\FloatTok{300}\NormalTok{))}

\CommentTok{\# Add axis for fitness}
\NormalTok{ax1 }\OperatorTok{=} \FunctionTok{Axis}\NormalTok{(}
\NormalTok{    fig[}\FloatTok{1}\NormalTok{, }\FloatTok{1}\NormalTok{],}
\NormalTok{    xlabel}\OperatorTok{=}\StringTok{"step"}\NormalTok{,}
\NormalTok{    ylabel}\OperatorTok{=}\StringTok{"fitness"}\NormalTok{,}
\NormalTok{    yticklabelsvisible}\OperatorTok{=}\ConstantTok{false}\NormalTok{,}
\NormalTok{)}

\CommentTok{\# Add axis for genotype{-}phenotype density}
\NormalTok{ax2 }\OperatorTok{=} \FunctionTok{Axis}\NormalTok{(}
\NormalTok{    fig[}\FloatTok{1}\NormalTok{, }\FloatTok{2}\NormalTok{],}
\NormalTok{    xlabel}\OperatorTok{=}\StringTok{"step"}\NormalTok{,}
\NormalTok{    ylabel}\OperatorTok{=}\StringTok{"genotype{-}phenotype density"}\NormalTok{,}
\NormalTok{    yticklabelsvisible}\OperatorTok{=}\ConstantTok{false}\NormalTok{,}
\NormalTok{)}

\CommentTok{\# Plot fitness}
\FunctionTok{lines!}\NormalTok{(ax1, F\_traj, color}\OperatorTok{=}\NormalTok{ColorSchemes.seaborn\_colorblind[}\FloatTok{1}\NormalTok{])}

\CommentTok{\# Plot genotype{-}phenotype density}
\FunctionTok{lines!}\NormalTok{(ax2, M\_traj, color}\OperatorTok{=}\NormalTok{ColorSchemes.seaborn\_colorblind[}\FloatTok{2}\NormalTok{])}

\NormalTok{fig}
\end{Highlighting}
\end{Shaded}

\includegraphics[width=12.5in,height=6.25in]{metropolis_kimura_evolution_files/figure-pdf/cell-11-output-1.png}

Let's now repeat the simulation multiple times and plot all different
trajectories.

\begin{Shaded}
\begin{Highlighting}[]
\BuiltInTok{Random}\NormalTok{.}\FunctionTok{seed!}\NormalTok{(}\FloatTok{42}\NormalTok{)}

\CommentTok{\# Define number of simulations}
\NormalTok{n\_sim }\OperatorTok{=} \FloatTok{5}

\CommentTok{\# Run simulations}
\NormalTok{x\_traj\_list }\OperatorTok{=}\NormalTok{ [}
    \FunctionTok{evo\_metropolis\_kimura}\NormalTok{(x0, fit\_peak, mut\_peaks, N, β, µ, n\_steps)}
\NormalTok{    for \_ }\KeywordTok{in} \FloatTok{1}\OperatorTok{:}\NormalTok{n\_sim}
\NormalTok{]}

\CommentTok{\# Initialize figure}
\NormalTok{fig }\OperatorTok{=} \FunctionTok{Figure}\NormalTok{(size}\OperatorTok{=}\NormalTok{(}\FloatTok{600}\NormalTok{, }\FloatTok{300}\NormalTok{))}

\CommentTok{\# Add axis for fitness landscape}
\NormalTok{ax1 }\OperatorTok{=} \FunctionTok{Axis}\NormalTok{(}
\NormalTok{    fig[}\FloatTok{1}\NormalTok{, }\FloatTok{1}\NormalTok{],}
\NormalTok{    xlabel}\OperatorTok{=}\StringTok{"phenotype 1"}\NormalTok{,}
\NormalTok{    ylabel}\OperatorTok{=}\StringTok{"phenotype 2"}\NormalTok{,}
\NormalTok{    aspect}\OperatorTok{=}\FunctionTok{AxisAspect}\NormalTok{(}\FloatTok{1}\NormalTok{),}
\NormalTok{    title}\OperatorTok{=}\StringTok{"Fitness"}\NormalTok{,}
\NormalTok{    yticklabelsvisible}\OperatorTok{=}\ConstantTok{false}\NormalTok{,}
\NormalTok{    xticklabelsvisible}\OperatorTok{=}\ConstantTok{false}\NormalTok{,}
\NormalTok{)}
\CommentTok{\# Add axis for genotype{-}phenotype density}
\NormalTok{ax2 }\OperatorTok{=} \FunctionTok{Axis}\NormalTok{(}
\NormalTok{    fig[}\FloatTok{1}\NormalTok{, }\FloatTok{2}\NormalTok{],}
\NormalTok{    xlabel}\OperatorTok{=}\StringTok{"phenotype 1"}\NormalTok{,}
\NormalTok{    ylabel}\OperatorTok{=}\StringTok{"phenotype 2"}\NormalTok{,}
\NormalTok{    aspect}\OperatorTok{=}\FunctionTok{AxisAspect}\NormalTok{(}\FloatTok{1}\NormalTok{),}
\NormalTok{    title}\OperatorTok{=}\StringTok{"Genotype{-}phenotype}\SpecialCharTok{\textbackslash{}n}\StringTok{Density"}\NormalTok{,}
\NormalTok{    yticklabelsvisible}\OperatorTok{=}\ConstantTok{false}\NormalTok{,}
\NormalTok{    xticklabelsvisible}\OperatorTok{=}\ConstantTok{false}\NormalTok{,}
\NormalTok{)}

\CommentTok{\# Plot fitness landscape}
\FunctionTok{heatmap!}\NormalTok{(ax1, x, y, F, colormap}\OperatorTok{=:}\NormalTok{algae)}
\CommentTok{\# Plot heatmap of genotype{-}phenotype density}
\FunctionTok{heatmap!}\NormalTok{(ax2, x, y, M, colormap}\OperatorTok{=}\FunctionTok{Reverse}\NormalTok{(ColorSchemes.Purples\_9))}

\CommentTok{\# Plot contour plot}
\FunctionTok{contour!}\NormalTok{(ax1, x, y, F, color}\OperatorTok{=:}\NormalTok{white)}
\FunctionTok{contour!}\NormalTok{(ax2, x, y, M, color}\OperatorTok{=:}\NormalTok{white)}

\CommentTok{\# Loop over simulations}
\ControlFlowTok{for}\NormalTok{ (i, x\_traj) }\KeywordTok{in} \FunctionTok{enumerate}\NormalTok{(x\_traj\_list)}
    \CommentTok{\# Plot trajectory}
    \FunctionTok{scatterlines!}\NormalTok{.(}
\NormalTok{        [ax1, ax2],}
        \FunctionTok{Ref}\NormalTok{(x\_traj[}\FloatTok{1}\NormalTok{, }\OperatorTok{:}\NormalTok{]),}
        \FunctionTok{Ref}\NormalTok{(x\_traj[}\FloatTok{2}\NormalTok{, }\OperatorTok{:}\NormalTok{]),}
\NormalTok{        color}\OperatorTok{=}\NormalTok{ColorSchemes.seaborn\_colorblind[i],}
\NormalTok{        markersize}\OperatorTok{=}\FloatTok{3}
\NormalTok{    )}
\ControlFlowTok{end}

\CommentTok{\# Set limits}
\FunctionTok{xlims!}\NormalTok{(ax1, }\OperatorTok{{-}}\FloatTok{4}\NormalTok{, }\FloatTok{4}\NormalTok{)}
\FunctionTok{ylims!}\NormalTok{(ax1, }\OperatorTok{{-}}\FloatTok{4}\NormalTok{, }\FloatTok{4}\NormalTok{)}
\FunctionTok{xlims!}\NormalTok{(ax2, }\OperatorTok{{-}}\FloatTok{4}\NormalTok{, }\FloatTok{4}\NormalTok{)}
\FunctionTok{ylims!}\NormalTok{(ax2, }\OperatorTok{{-}}\FloatTok{4}\NormalTok{, }\FloatTok{4}\NormalTok{)}

\NormalTok{fig}
\end{Highlighting}
\end{Shaded}

\includegraphics[width=12.5in,height=6.25in]{metropolis_kimura_evolution_files/figure-pdf/cell-12-output-1.png}

\subsection{Conclusion}\label{conclusion}

In this notebook, we developed a biologically-motivated model for
evolutionary dynamics that combines elements of the Metropolis-Hastings
algorithm with classic population genetics theory. This model implements
a two-step process that separately accounts for:

\begin{enumerate}
\def\labelenumi{\arabic{enumi}.}
\tightlist
\item
  The probability of mutation occurring (based on genotype-phenotype
  density)
\item
  The probability of fixation within a population (based on Kimura's
  fixation probability)
\end{enumerate}

By separating these processes, our \emph{Metropolis-Kimura} algorithm
provides a plausible model of evolutionary dynamics, where mutational
constraints and selection pressures interact to shape evolutionary
trajectories. This approach allows us to explore how both the
accessibility of phenotypic space and the fitness landscape influence
the paths taken by evolving populations.



\end{document}
